\section{接線型の第二共鳴と一般点で二重な成分}
\label{sec:tangent-double-v34}

この節では重みを $(1,1,1,3)$ に固定する。
Poisson スキームを $\mathcal P$ とし、その座標と混合式を
\[
 (c,D,[F])\in S_3\oplus\mathfrak{sl}_3\oplus(S_4^{\oplus3}/xS_3),
 \qquad M(c,D;F)=2F(c)-(x\times D)\cdot F-c\operatorname{div}F
\]
と書く。方程式は $D(c)=M(c,D;F)=0$ であり、イデアルは根基化しない。
前節で残した接線型の第二の共鳴 $a=-3/2$ は、
既知成分の内部の核の増大だけではなく、別の非被約成分を与える。

\begin{theorem}[一般点で二重な十八次元成分]
\label{thm:tangent-double-v34}
\[
 q=xz-y^2,\quad c_0=xq,\quad N_0=x\partial_y+2y\partial_z,\quad
 D_-=-\tfrac32N_0,\quad F_-=(0,\tfrac12y^3z,y^2z^2)
\]
と置く。組 $(c_0,D_-)$ の $\operatorname{SL}_3$ 軌道上で全混合核を動かした
十八次元族の閉包 $\mathcal E_{-3/2}$ は、$\mathcal P$ の既約成分である。
その一般点の局所環は
\[
 K[\epsilon]/(\epsilon^2)
\]
であり、一般点の長さは二である。被約成分は重み付き座標変換群の稠密軌道を持つ。

より正確には、$P_-=(c_0,D_-,0)$ における全完備局所環は
\begin{equation}
 \widehat{\mathcal O}_{\mathcal P,P_-}
 \simeq
 \mathbb C[[s_1,\ldots,s_{17},k,b,v]]/(kv,b^2v).
 \label{eq:tangent-double-origin-v34}
\end{equation}
$v=0$ の枝は十九次元成分 $L(1,2,3)$ であり、もう一方は
$\mathcal E_{-3/2}$ の二重構造である。
点 $Q_-=(c_0,D_-,F_-)$ では
\begin{equation}
 \widehat{\mathcal O}_{\mathcal P,Q_-}
 \simeq\mathbb C[[s_1,\ldots,s_{18},b]]/(b^2).
 \label{eq:tangent-double-nonzero-v34}
\end{equation}
従ってこの成分は、既知の九成分と相異なる。
重み三の Poisson スキームには少なくとも十個の既約成分がある。
\end{theorem}

この定理での長さ二は、混合行列の一点での階数だけから推測したものではない。
以下の消去では、二つの基底変数を含む全形式近傍の式を求める。

\subsection{非零の接線型に共通する基底の形式切断}

\begin{lemma}
\label{lem:tangent-base-slice-v34}
$A\ne0$ を固定する。基底 $D(c)=0$ の $(c_0,AN_0)$ における形式近傍は、
七個の滑らかな軌道変数と、次の二変数族の積である。
\begin{equation}
 q_b=xz-y^2+\frac{b}{2A}z^2,\qquad
 c=u xq_b,\qquad D=(0,Ax+bz,2Ay),\qquad (u,b)=(1,0).
 \label{eq:tangent-base-slice-v34}
\end{equation}
特に基底はこの点で滑らかな九次元スキームである。
\end{lemma}

\begin{proof}
$D$ は正則冪零行列なので、共役作用の六変数を除き、形式的に
\[
 D=(tz,Ax+bz,2Ay)
\]
と書ける。実際、その二つの行列不変量は
$\operatorname{tr}(D^2)=4Ab$、$\det D=2A^2t$ であり、
正則行列の共役軌道への横断切断を与える。

三次係数を単項式 $x^iy^jz^{3-i-j}$ の順に並べる。
$x^2z,x^3$ の係数を $u,r$ として残し、他の八係数を消去する。
使用する八次小行列式は $b=t=0$ で $-1152A^8$ なので単元である。
消去後の三次式を $4A^2$ 倍すると
\begin{align*}
 4A^2c={}&4A^2(ux^2z-uxy^2+rx^3)+2Abu xz^2
       +2t^2r z^3+t(Au-6br)yz^2\\
       &+8Atr y^3-12Atr xyz.
\end{align*}
これに $D$ を作用させた残差は
\[
 3bt(Au-2br)z^3+9At(Au-2br)xz^2
\]
のみである。$Au-2br$ は対象の点で単元なので、残る条件は正確に $t=0$ となる。
従って $c=uxq_b+rx^3$ である。

$D$ の中心化代数の場 $D^2-\operatorname{tr}(D^2)I/3$ は、
$b=0$ で $c$ の $x^3$ 係数を $2A^2u$ だけ動かす。
この係数は単元なので、中心化群の一変数によって $r$ を形式的に零にできる。
以上で七軌道変数と式 \eqref{eq:tangent-base-slice-v34} を得る。
独立の確認として、$D(c)$ の九次小行列式は $2592A^8$、
七軌道方向と $\partial_u,\partial_b$ の接写像の九次小行列式は $-6$ となる。
従ってこの切断は被約点だけの標準形ではなく、全形式近傍を記述している。
\end{proof}

$b$ は幾何的には直線と二次曲線の接触を解く変数である。
$b=0$ では $x=0$ と $q_b=0$ は接し、$b\ne0$ では二点で横断的に交わる。
また $u$ は三次式と線形場との相対的な係数を記録する。

\subsection{二重構造を決める正確な混合方程式}

補題の $A=-3/2$ を取る。二次場の係数は
\[
 p=\frac{\nabla c-x\times D}{3}
\]
であり、三次式による $w$ 移動の混合項
$p\times\nabla g$、$g\in S_3$ は、全 $(u,b)$ 上で実際の核を成す。
その十方向は $(1,0)$ で独立なので、形式近傍で直和因子として除ける。
残る二十五列に対し、二十四の単元行・列を消去する。
最後の一方向には次の有理係数の四次場を取れる。
\begin{align}
 \widehat F_1={}&0,\nonumber\\
 \widehat F_2={}&\frac{2u}{u+3}y^3z
 -\frac{(u-1)(4u+3)}{2(u+1)(u+3)}xyz^2
 -\frac{b(7u^2+6u-9)}{3(u+1)(u+3)(5u-3)}yz^3,\label{eq:tangent-double-Fhat-v34}\\
 \widehat F_3={}&y^2z^2+\frac{bu(9-5u)}{3(u+3)(5u-3)}z^4.\nonumber
\end{align}
全ての分母は $(u,b)=(1,0)$ で単元であり、同点で $\widehat F=F_-$ となる。
直接の係数比較により、正確な恒等式
\begin{align}
 M(c,D;\widehat F)
 ={}&\frac{2b^2u(u-3)(4u+3)}{9(u+1)(u+3)(5u-3)}xz^5\nonumber\\
 &+\frac{u(u-1)(4u+3)}{2(u+1)(u+3)}x^3z^3
 \label{eq:tangent-double-Schur-v34}
\end{align}
を得る。

この式が Schur 消去の全残差であることも、小行列式で保証できる。
Euler 商の基底は成分順、単項式順とし、第一成分の $x$ で割れる単項式を除く。
零始まりで、単元列は
\[
 \{0,1,\ldots,20,24,31,34\}
\]
であり、単元行は $S_6$ の単項式順の
\[
 \{2,3,4,5,6,8,9,10,11,12,13,14,15,16,17,19,20,21,22,23,24,25,26,27\}
\]
である。この二十四次小行列式は $-136193423625$。
十個の $w$ 移動、これらの二十四列、$\widehat F$ の三十五次基底変換の行列式は
$-16875$ である。式 \eqref{eq:tangent-double-Schur-v34} の二行は上の単元行に含まれず、
その二行以外の係数は恒等的に零である。
従ってこの基底変換と行消去は形式近傍で可逆であり、隠れた高次方程式は残らない。

$\widehat F$ の係数を $v$、$k=u-1$ とする。
$b^2$ と $k$ の係数は原点でそれぞれ $-7/36,7/16$ なので単元である。
従って混合条件は正確に
\[
 b^2v=kv=0
\]
へ変形できる。七軌道変数と十個の $w$ 移動を戻せば
式 \eqref{eq:tangent-double-origin-v34} を得る。
ここで全ての基底変換と単元小行列式は $(u,b)$ だけに依存し、
消去は $F$ に関して線形である。従って同じ操作は基底の形式近傍と全混合空間の積上で成立し、
$v=0$ だけでなく $v=1$ の近傍にも適用できる。

\begin{proof}[定理 \ref{thm:tangent-double-v34} の残りの証明]
$v\ne0$ の点では $k=0,b^2=0$ となるので、
式 \eqref{eq:tangent-double-nonzero-v34} が従う。
$(c_0,D_-)$ の軌道は七次元で、混合核の階数は十一。
従ってその核束の全体は既約十八次元族である。
上の局所環から、その閉包は全スキームの一つの既約成分であり、一般点の長さが二と分かる。

混合核は十個の $w$ 移動と $F_-$ で生成される。
残余係数が非零なら、移動によって十方向を除き、$w$ の非零定数倍で残余係数を一にできる。
従って被約成分には稠密な座標変換軌道がある。
従来の八成分は次元が十九以上であり、既知の接線型十八次元成分
$\mathcal E_{-5/2}$ は一般点で滑らかである。
いずれも一般点で二重な $\mathcal E_{-3/2}$ と一致しない。

最後に $v=0$ の枝を同定する。$b\ne0$ では直線・二次曲線が横断し、
$u$ は $-3/2$ 近傍の非共鳴な比を記録するため、混合核は全て $w$ 移動から来る。
これは十九次元対数型成分 $L(1,2,3)$ の稠密部分である。
よってその閉包が式 \eqref{eq:tangent-double-origin-v34} の $v=0$ の枝である。
\end{proof}

さらに
\[
 (kv,b^2v)=(v)\cap(k,b^2)
\]
は primary 分解である。二枝の被約交叉は滑らかな十七次元で、primary 成分同士の交叉は
その上の $\mathbb C[[b]]/(b^2)$ という二重構造を持つ。
また $Q_-$ で全 Jacobi 行列の階数は三十四、接空間次元は十九であり、
三十四次小行列式 $27799801630335/64$ が独立の確認を与える。

\paragraph{二次で止まる接触解消方向.}
$u=1$ における式 \eqref{eq:tangent-double-Fhat-v34} は単純化し、
\[
 c_b=x(xz-y^2-bz^2/3),\quad
 D_b=(0,-3x/2+bz,-3y),\quad
 F_b=(0,y^3z/2-byz^3/12,y^2z^2+bz^4/6)
\]
となる。この族は正確に
\[
 D_b(c_b)=0,\qquad M(c_b,D_b;F_b)=-\frac7{36}b^2xz^5
\]
を満たす。従って接触を解く一次方向は存在するが、この点では二次で止まる。
この障害を消して被約化してしまうと、新成分の長さ二は見えなくなる。

\subsection{他の接線共鳴との比較}

同じ基底の切断は、他の非零 $A$ にも適用できる。
$P_A=(c_0,AN_0,0)$ での全形式局所環を、次の表にまとめる。
$k=u-1$ であり、$s_1,\ldots,s_{17}$ は七軌道変数と十移動変数である。
\[
\begin{array}{c|c|c}
 A&\text{滑らかな十七変数を除いた局所環}&\text{通る大域成分}\\\hline
 2&\mathbb C[[b,k,v_1,\ldots,v_{10}]]/(kv_1,\ldots,kv_{10})
   &L(1,2,3),\ R(1,5)\\
 -1&\mathbb C[[b,k,v_1,\ldots,v_4]]/(kv_1,\ldots,kv_4)
   &L(1,2,3),\ R(2,4)\\
 -5/2&\mathbb C[[b,k,v]]/(kv,bv)&L(1,2,3),\ \mathcal E_{-5/2}\\
 -3/2&\mathbb C[[b,k,v]]/(kv,b^2v)&L(1,2,3),\ \mathcal E_{-3/2}
\end{array}
\]

$A=-5/2$ の場合には
\[
 \widehat F=(0,4uyz^3/(3u+5),z^4)
\]
を取れば、全残差は
\[
 M(c,D;\widehat F)
 =\frac{b(u-5)(4u+5)}{5(3u+5)}xz^5
  -\frac{5(u-1)(4u+5)}{2(3u+5)}x^2z^4
\]
となる。二つの係数は $b,k$ の単元倍なので、こちらの十八次元成分は被約である。

$A=2,-1$ については、それぞれ
\[
 \nabla x\times\nabla S_5,\qquad \nabla q_b\times\nabla S_4
\]
が $u=1$ 上の全混合核を与える。
十移動方向を除く残余次元は十、四である。
これらの残余方向を全形式基底へ延長したとき、混合像は全て $u-1$ で割れる。
その商を単元像で割った行列は、原点でそれぞれ階数十、四を持つ。
従って追加の基本因子は全て一位の $k$ であり、表の被約な局所環を得る。
この階数にはそれぞれ二十五次小行列式
$-8321693751705600$、$-5015306502144$ を用いる。

$A\ne0,2,-1,-5/2,-3/2$ では、既に求めた全パラメータの重みブロック表により
混合行列の階数は二十五である。基底切断の一般点では十個の独立な移動方向があるので、
二十六次小行列式は恒等的に消える。一方、対象点の二十五次小行列式は単元である。
従って混合行列の階数はこの基底近傍で一定で、核は階数十のベクトル束となる。
その一般点は $L(1,2,3)$ に属すので、$P_A$ は同成分上の滑らかな十九次元の点である。
これは有限個の非共鳴値を試す議論ではない。

この説明では $A=1$ を含めている。この値では $p=(q,0,0)$ なので
$p\times\nabla x^3=0$ となり、実際の移動の像は九次元である。
不足する核方向は $x^3N_0$ であり、
\[
 p_a\times\nabla x^3=(1-a)x^3N_0
\]
によって近傍の移動方向の有限な極限として回収できる。
従って移動の階数と混合核の次元をこの点で混同してはならない。

$A=0$ では基底補題は用いないが、固定した $(c_0,0)$ 上の混合核は全て
$L(1,2,3)$ の閉包に含まれることを明示できる。
この核は九移動方向と $c_0N_0$ で生成され、
\[
 p_a\times\nabla c_0=-a c_0N_0,\qquad
 p_a=(\nabla c_0-a x\times N_0)/3
\]
である。よって $a\ne0$ での移動 $g=-c_0/a$ の混合項は $c_0N_0$ であり、
$a\to0$ の有限な極限が残る方向を覆う。
ここでは $D=0$ における全局所環の分類までは主張しない。

\subsection{二次元 Lie 作用と有理第一積分}

\begin{proposition}
点 $Q_-$ の発散零持ち上げは
\[
 \widehat P_-=H\wedge N,\qquad
 H=\tfrac1{12}(7x\partial_x+y\partial_y-5z\partial_z-9w\partial_w),
 \qquad N=x\partial_y+2y\partial_z-y^2z\partial_w
\]
である。$[H,N]=N/2$、$\operatorname{div}H=-1/2$、$\operatorname{div}N=0$ を満たす。
また
\[
 J=x^2w+\frac13xy^3z-\frac2{15}y^5
   =x^2w+\frac13q y^3+\frac15y^5
\]
に対し $N(q)=N(J)=0$、$H(q)=q/6$、$H(J)=5J/12$ である。
対応する水平一形式 $\alpha$ は
\[
 2q\,dJ-5J\,dq=2x\alpha
\]
を満たし、$J^2/q^5$ が有理第一積分となる。
\end{proposition}

\begin{proof}
持ち上げの六係数、Lie 括弧、発散、および一形式の四係数を直接比較すればよい。
$N$ は補助次数 $(0,1,2,5)$ を一つ下げるため局所冪零である。
\end{proof}

これは次数 $(2,5)$ の特殊 pencil から共通因子 $x$ を除いた次数六の一形式であり、
通常の $R(2,5)$ と同一視しない。
重み付き座標変更 $w'=-5w-yz^2$ は $H$ と可換であり、
$N(w')=-qz$ へ変える。この表示は、同じ成分の量子化と不変式環を調べる際にも使える。

\paragraph{厳密検証と射程.}
\texttt{verify\_tangent\_resonance\_v34.py} は Python 標準ライブラリのみで動く。
分母を払った全係数恒等式、基底変換と単元部分の小行列式、四共鳴の局所式、
全 Jacobi 行列、Lie 表示と第一積分を有理数上で確認する。
通常実行では入力や保存済み証明書を書き換えない。
本節は接線型の非零線形場の形式近傍を決定するものであり、
全ての非被約三次式の層や全 Poisson 成分の完全分類を主張するものではない。
